\section{Sketch Comparisons and Classical Controls}
\label{app:v8-sketch-comparisons}

The sketch experiments are controls for the state-level story. They are not
the main LLM application. Their role is to show the difference between
state-level mergeability, scalar-output collapse, and learned approximation.

\subsection{Generated Sketch Grid}

The generated sketch assets live under \texttt{assets/sketches/}.
Table~\ref{tab:v8-sketch-compact} is the compact grid; the figures break out
detail by query family.

\begin{table}[t]
\centering
\resizebox{\linewidth}{!}{% Auto-generated by treepo.bench.reports.classical_sketches; do not edit.
\begin{tabular}{lllrrrrrrrr}
\toprule
family & query & best official & official & learned $f$ & $f$ RMSE & learned $g$ & $g$ RMSE & learned joint (best) & joint variant & joint RMSE \\
\midrule
distinct & cardinality & theta\_datasketches & 0 & -- & -- & learned\_g\_exact\_distinct\_union\_state\_space & 0 & learned\_joint\_exact\_distinct\_union\_state\_space & fg & 0 \\
frequency & focus\_frequency & -- & -- & -- & -- & learned\_g\_count\_min\_state\_space & 0 & learned\_joint\_count\_min\_state\_space & fg & 0 \\
frequency & top5\_point\_frequency & count\_min\_datasketches & 0 & -- & -- & -- & -- & -- & -- & -- \\
quantile & rank\_at\_q0.5 & quantiles\_floats\_datasketches & 0.002583 & -- & -- & -- & -- & learned\_joint\_quantiles\_reference\_q0.5 & fg & 0.02422 \\
quantile & rank\_at\_q0.95 & tdigest\_double\_datasketches & 0.002462 & -- & -- & -- & -- & learned\_joint\_kll\_reference\_q0.95 & fg & 0.007558 \\
sampling & accumulator\_summary\_sum & tuple\_accumulator\_datasketches & 0 & -- & -- & -- & -- & learned\_joint\_tuple\_summary\_sum\_reference & fg & 0.007204 \\
sampling & total\_weight & varopt\_strings\_datasketches & 0 & -- & -- & learned\_g\_exact\_total\_weight\_state\_space & 0 & learned\_joint\_exact\_total\_weight\_state\_space & fg & 0 \\
set & a\_not\_b & theta\_datasketches & 0 & -- & -- & -- & -- & learned\_joint\_exact\_set\_a\_not\_b & fg & 0.09717 \\
set & intersection & theta\_datasketches & 0 & -- & -- & -- & -- & learned\_joint\_exact\_set\_intersection & fg & 0.2417 \\
set & union & theta\_datasketches & 0 & -- & -- & -- & -- & learned\_joint\_exact\_set\_union & fg & 0.03052 \\
\bottomrule
\end{tabular}
}
\caption{\textbf{Generated compact sketch grid.} Official rows are reference
sketch reductions or controls. Learned rows test whether the selected state
and query class can recover the same behavior under tree reductions.}
\label{tab:v8-sketch-compact}
\end{table}

\begin{figure}[t]
\centering
\includegraphics[width=0.48\linewidth]{classical_sketches_distinct_detail.pdf}\hfill
\includegraphics[width=0.48\linewidth]{classical_sketches_frequency_detail.pdf}
\caption{\textbf{Distinct-count and frequency sketch controls.} The plots
separate state-level parity from scalar output recovery.}
\label{fig:v8-sketch-distinct-frequency}
\end{figure}

\begin{figure}[t]
\centering
\includegraphics[width=0.48\linewidth]{classical_sketches_quantile_detail.pdf}\hfill
\includegraphics[width=0.48\linewidth]{classical_sketches_set_detail.pdf}
\caption{\textbf{Quantile and set-style sketch controls.} Randomized and
rank-preserving families illustrate where tree reductions need a state, not a
scalar answer.}
\label{fig:v8-sketch-quantile-set}
\end{figure}

\begin{figure}[t]
\centering
\includegraphics[width=0.64\linewidth]{classical_sketches_sampling_detail.pdf}
\caption{\textbf{Sampling-style sketch controls.} The diagnostic emphasizes
the design distinction between query correctness and sampling variability.}
\label{fig:v8-sketch-sampling}
\end{figure}

\subsection{Reading the Comparisons}

An exact classical row means the state, merge, and readout are supplied by the
reference sketch. A learned row means the experiment searches within a selected
state/query class for a tree-compatible approximation. If the learned row fails
to match the reference, the failure may be representational, optimization-side,
or caused by the query shape itself. The C-TreePO certificate handles the
deployed tree by auditing its realized local-law residuals.
